\scanpage{289}
\mspage{205}
% archive sheet number: 288
4) Considérons $M=\mathrm{Ind}_{\unsure{\pi|\kappa|k}_0}(C)$, \deleted{\unsure{$C$ Karoubi}} \inserted{$C$ petite,} et considérons la filtration canonique $(M_\pi)_{\pi\geq\pi_0}$. \deleted{\unsure{Y a-t}} Donner des exemples \deleted{de cardinaux} où il y a des cardinaux non utiles, ou mieux, où pour \unsure{tt|$\forall$} cardinal $\pi$, existe un cardinal $\pi'>\pi$ qui soit non utile. (Cela donnera l'exemple demandé dans $1^\circ$.) D'autre part, donner des exemples de cardinaux utiles qui ne soient réguliers -- si possible, comme tantôt, \deleted{des} tels cardinaux aussi grands qu'on veut. \par\hfill\unsure{$\backslash$}\hfill\null

\uline{Définition} 4.8.\unsure{?|7}. Soit $M$ une $\unsure{\mathcal U|U}$-catégorie, \inserted{$\pi_0$ un cardinal infini.} On dit que $M$ est $\pi_0$-\uline{accessible} si $M$ est équivalente à une catégorie $\mathrm{Ind}_{\pi_0}(C)$, \unsure{-|où} une petite catégorie, i.e si elle satisfait aux conditions équivalentes b, c, c' de 4.8.1 (impliquant la condition a), sans doute strictement plus faible). On dit que $M$ est \uline{accessible} s'il existe un cardinal tel que $M$ \unsure{soit|est} $\pi_0$-accessible [NB si $\pi_0$ est un tel cardinal, alors tout cardinal $\pi>\pi_0$ tel que $\pi^{\pi_0}=\pi$, notamment les cardinaux $\pi'^{\pi_0}$ avec $\pi'\geq2$, ont la $=$ propriété.], ou ce
